\subsection{Уровень риска: сколько стоит гарантия}\label{sec:alpha}

Миссии и покрытие пересчитаны при $\alpha \in \{0{,}05;\,0{,}1;\,0{,}2;\,0{,}3\}$ (рис.~\ref{fig:alpha}, табл.~\ref{tab:alpha}; дрейф L3). При $\alpha = 0{,}05$ нужно не меньше 19 калибровочных сцен, поэтому разбиение 19/2, при остальных --- 13/8.

\begin{figure}[H]
\centering
\includegraphics[width=\linewidth]{figures/alpha_sweep.png}
\caption{Уровень риска $\alpha$ против покрытия (слева, тонкая линия --- номинал $1-\alpha$), радиуса запаса (в центре) и исхода миссий (справа); дрейф L3, растение (сплошная) и велосипед (пунктир)}
\label{fig:alpha}
\end{figure}

\begin{table}[H]
\centering\small
\caption{Цена гарантии при разных $\alpha$ (L3): покрытие и радиус совместной ветки, доля отказов по растению, доля задач с найденным путём и доля нарушений}
\label{tab:alpha}
\begin{tabular}{c|cc|cc|c|cc|cc}
\toprule
$\alpha$ & \multicolumn{2}{c|}{покрытие} & \multicolumn{2}{c|}{$\hat r$, м} & отказ & \multicolumn{2}{c|}{трудные: путь} & \multicolumn{2}{c}{случайные: успех} \\
 & раст. & велос. & раст. & велос. & раст. & совм. & оракул & совм. & оракул \\
\midrule
0{,}05 & 0{,}95 & 0{,}95 & 0{,}85 & 0{,}76 & 0{,}70 & 0{,}10 & 0{,}99 & 0{,}34 & 0{,}97 \\
0{,}10 & 0{,}93 & 0{,}93 & 0{,}80 & 0{,}75 & 0{,}52 & 0{,}30 & 0{,}99 & 0{,}68 & 0{,}97 \\
0{,}20 & 0{,}86 & 0{,}86 & 0{,}74 & 0{,}70 & 0{,}00 & 0{,}31 & 0{,}99 & --- & --- \\
0{,}30 & 0{,}72 & 0{,}72 & 0{,}60 & 0{,}50 & 0{,}00 & 0{,}31 & 0{,}99 & 0{,}71 & 0{,}97 \\
\bottomrule
\end{tabular}

\end{table}

\begin{itemize}[nosep]
  \item Совместная ветка держит номинал при каждом $\alpha$ (0{,}95; 0{,}94; 0{,}86--0{,}87; 0{,}72), раздельная остаётся консервативной и при $\alpha = 0{,}3$ (0{,}82--0{,}91).
  \item Отказы по растению исчезают при $\alpha \ge 0{,}2$: сцена, где детектор пропускает растение целиком, перестаёт определять квантиль.
  \item Радиус убывает медленно: от 0{,}81 до 0{,}60~м для растения и от 0{,}75 до 0{,}50~м для велосипеда при росте $\alpha$ с 0{,}1 до 0{,}3. Решаемость миссий почти не растёт: 30~$\to$~31\,\% задач. Ниже $\alpha = 0{,}1$ она резко падает: 10\,\% при $\alpha = 0{,}05$.
  \item Нарушения у совместной ветки остаются на уровне $0{,}2\,\%$ задач даже при $\alpha = 0{,}3$, потому что гарантия по классу и сцене консервативна по отношению к отдельной задаче.
\end{itemize}
Вывод: торговаться уровнем риска почти бесполезно. Решаемость упирается в «пол» ошибки восприятия: даже при $\alpha = 0{,}3$ истинный след выходит за область класса на 0{,}5--0{,}6~м. Сократить запас может только лучшее восприятие, а не более смелая калибровка.

\subsection{Стресс-тест композиции}\label{sec:composition}

На синтетическом дрейфе раздельная калибровка переплачивает радиусом, но не недопокрывает (разд.~\ref{sec:coverage}). Причина в том, что дрейф независим от ошибок детектора. Чтобы проверить, когда сумма квантилей перестаёт быть валидной, ошибки меток оставлены реальными, а позе задана сценовая смесь. В $k$ сценах из 21 происходит «отказ локализации»: реализации берутся с уровня L5. В остальных сценах --- с уровня L1. Зависимость задаётся выбором этих сцен:
\begin{itemize}[nosep]
  \item \emph{согласованно}: там, где метки хуже всего;
  \item \emph{разнесённо}: там, где метки лучше всего;
  \item \emph{случайно}: случайное подмножество.
\end{itemize}
Уровень риска $\alpha = 0{,}3$ взят потому, что при $\alpha = 0{,}1$ и 13 сценах квантиль равен максимуму и редкие отказы всегда попадают в радиус.

\begin{figure}[H]
\centering
\includegraphics[width=\linewidth]{figures/composition_stress_a03_L5.png}
\caption{Покрытие раздельной (розовая) и совместной (синяя) калибровки в зависимости от доли сцен с отказом локализации и от того, в каких сценах происходят отказы; $\alpha = 0{,}3$, нижняя пунктирная линия --- граница $1-2\alpha$, которую гарантирует сумма квантилей}
\label{fig:composition}
\end{figure}

Покрытие раздельной схемы зависит от того, как распределены отказы (рис.~\ref{fig:composition}):
\begin{itemize}[nosep]
  \item при согласованных отказах оно 0{,}80--0{,}84;
  \item при разнесённых падает до 0{,}66 для растения (номинал 0{,}70) при 19--24\,\% сцен с отказом.
\end{itemize}
Совместная ветка держит 0{,}71--0{,}78 во всех конфигурациях. Теоретически сумма двух квантилей гарантирует только $1 - 2\alpha$: если промах совместной карты не больше суммы промахов меток и позы, то по неравенству Буля $P(\text{оба не превышены}) \ge 1 - 2\alpha$. Достигается ли эта граница, решает зависимость ошибок, которую на практике заранее не знают. Совместная калибровка от неё не зависит. Это согласуется с выводом~\cite{baral} о том, что маргинальная калибровка модулей не переносится на систему. Наш случай --- недопокрытие при \emph{разнесённых} редких отказах.

\subsection{Проверка упрощений: маски, реальная одометрия, другие сцены}\label{sec:variants}

Основной эксперимент опирается на три упрощения: маски из рамок, синтетический дрейф и одно семейство сцен. Каждое из них проверено отдельно (табл.~\ref{tab:variants}).

\begin{table}[H]
\centering\scriptsize\setlength{\tabcolsep}{3.5pt}
\caption{Проверка упрощений: покрытие и медианный радиус (м) для велосипеда и растения, $\alpha = 0{,}1$; $^\dagger$ --- ветка отказывается в половине разбиений или чаще, приведён радиус в остальных}
\label{tab:variants}
\begin{tabular}{l|cccccc|cccccc}
\toprule
Вариант & \multicolumn{6}{c|}{велосипед} & \multicolumn{6}{c}{растение} \\
 & \multicolumn{2}{c}{метки} & \multicolumn{2}{c}{раздельно} & \multicolumn{2}{c}{\textbf{совместно}} & \multicolumn{2}{c}{метки} & \multicolumn{2}{c}{раздельно} & \multicolumn{2}{c}{\textbf{совместно}} \\
 & покр. & $\hat q$ & покр. & $\hat q$ & покр. & $\hat q$ & покр. & $\hat q$ & покр. & $\hat q$ & покр. & $\hat q$ \\
\midrule
рамка + глубина, L0 & 0{,}93 & 0{,}75 & 0{,}93 & 0{,}75 & 0{,}93 & 0{,}75 & 0{,}93 & 0{,}85$^\dagger$ & 0{,}93 & 0{,}85$^\dagger$ & 0{,}93 & 0{,}85$^\dagger$ \\
рамка + глубина, L3 & 0{,}93 & 0{,}75 & 0{,}99 & 0{,}90 & 0{,}94 & 0{,}75 & 0{,}94 & 0{,}85$^\dagger$ & 0{,}96 & 1{,}09$^\dagger$ & 0{,}94 & 0{,}81$^\dagger$ \\
рамка + глубина, L4 & 0{,}76 & 0{,}75 & 0{,}97 & 1{,}61 & 0{,}93 & 1{,}18 & 0{,}86 & 0{,}85$^\dagger$ & 0{,}95 & 2{,}11$^\dagger$ & 0{,}93 & 1{,}53$^\dagger$ \\
MobileSAM, L0 & 0{,}93 & 0{,}75 & 0{,}93 & 0{,}75 & 0{,}93 & 0{,}75 & 0{,}93 & 0{,}61$^\dagger$ & 0{,}93 & 0{,}61$^\dagger$ & 0{,}93 & 0{,}61$^\dagger$ \\
MobileSAM, L3 & 0{,}93 & 0{,}75 & 0{,}97 & 0{,}90 & 0{,}93 & 0{,}75 & 0{,}94 & 0{,}61$^\dagger$ & 0{,}95 & 0{,}86$^\dagger$ & 0{,}95 & 0{,}61$^\dagger$ \\
MobileSAM, L4 & 0{,}78 & 0{,}75 & 0{,}98 & 1{,}61 & 0{,}93 & 1{,}03 & 0{,}82 & 0{,}61$^\dagger$ & 0{,}94 & 1{,}92$^\dagger$ & 0{,}95 & 1{,}63$^\dagger$ \\
ВО, каждый 2-й кадр & 0{,}67 & 0{,}75 & 0{,}92 & 2{,}46$^\dagger$ & 0{,}93 & 1{,}95$^\dagger$ & 0{,}74 & 0{,}85$^\dagger$ & 0{,}97 & 2{,}98$^\dagger$ & 0{,}96 & 2{,}31$^\dagger$ \\
ВО, каждый 4-й кадр & 0{,}51 & 0{,}75 & 0{,}97 & 3{,}15$^\dagger$ & 0{,}96 & 2{,}41$^\dagger$ & 0{,}75 & 0{,}85$^\dagger$ & 0{,}99 & 3{,}46$^\dagger$ & 0{,}97 & 2{,}42$^\dagger$ \\
HM3D (OOD), L0 & --- & --- & --- & --- & --- & --- & 0{,}50 & 0{,}85 & 0{,}50 & 0{,}85 & 0{,}50 & 0{,}85 \\
HM3D (OOD), L3 & --- & --- & --- & --- & --- & --- & 0{,}50 & 0{,}85 & 0{,}50 & 1{,}07 & 0{,}50 & 0{,}82 \\
\bottomrule
\end{tabular}

\end{table}

\textbf{Маски MobileSAM.} Рамки детектора переведены в маски MobileSAM~\cite{mobilesam} --- облегчённым SAM, как в ConceptGraphs.
\begin{itemize}[nosep]
  \item Типичный промах велосипеда без дрейфа уменьшился почти в 9 раз: медиана 0{,}43~$\to$~0{,}05~м. Наихудший не изменился (0{,}75~м). Для растения медиана 0{,}49~$\to$~0{,}32~м.
  \item При 13 калибровочных сценах и $\alpha = 0{,}1$ квантиль равен максимуму. Поэтому радиус велосипеда не изменился (0{,}75~м), а у растения уменьшился с 0{,}85 до 0{,}61~м.
  \item Покрытие совместной ветки сохранилось (0{,}93--0{,}95).
\end{itemize}
Значит, при малом числе калибровочных сцен запас определяет наихудшая сцена: улучшать нужно хвост ошибок восприятия, а не среднее.

Миссии с масками SAM (те же 233 задачи) показывают обе стороны.
\begin{itemize}[nosep]
  \item Без дрейфа карта без калибровки стала безопасной: 0 нарушений против 27\,\% с масками из рамок.
  \item При дрейфе нарушения возвращаются: 9{,}5\,\% найденных путей на L3 и 19{,}5\,\% на L4.
  \item Откалиброванные ветки по-прежнему находят путь лишь в 21--38\,\% задач, потому что радиус задаёт худшая сцена. Нарушения у них --- 1--2\,\% путей на L1--L3, в пределах гарантии по классу и сцене.
\end{itemize}
Хорошая сегментация снимает ошибку меток в типичной сцене, но не ошибку позы. Поэтому при дрейфе гарантия нужна и с хорошим сегментатором, а её консервативность при 13 сценах --- плата за малую калибровочную выборку.

\textbf{Реальная визуальная одометрия.} Синтетический дрейф заменён на RGB-D одометрию Open3D~\cite{open3d}: плотная фотометрическая и геометрическая невязка~\cite{steinbrucker}, параметры по умолчанию. Одометрия работает кадр к кадру, без замыканий циклов, и применяется к тем же сценам.
\begin{itemize}[nosep]
  \item Каждый 8-й кадр не годится: повороты на месте по $16^\circ$ за шаг срывают сопоставление, и ATE на dev-сцене доходит до 12~м. Поэтому скачан каждый 2-й кадр (RGB и глубина).
  \item Добавлено плоское ограничение наземного робота: каждое относительное движение проецируется на рысканье и сдвиг по полу. Без него одна неудачная пара кадров у стены опрокидывала траекторию на $89^\circ$.
  \item Параметры по истине \emph{не} подбирались: перебор на трудных участках не дал настройки, лучшей везде, а выбор по эталонным позам был бы подгонкой под тест.
\end{itemize}

Ошибка такой одометрии --- тяжёлый хвост по сценам: ATE от 3{,}5~см до 3{,}9~м, медиана 78~см, дрейф в конце 1--27\,\% пути. Ошибка случается в основном на отдельных шагах --- поворотах на месте вблизи стен. Вариант с каждым 4-м кадром даёт медиану 2{,}6~м и служит режимом «потерянной локализации». Результаты калибровки (табл.~\ref{tab:variants}, рис.~\ref{fig:variants}):
\begin{itemize}[nosep]
  \item ветка «только метки» \emph{молча} теряет покрытие: 0{,}67 для велосипеда и 0{,}74 для растения (ВО, каждый 2-й кадр), 0{,}51 для велосипеда (каждый 4-й). Это гипотеза H1 уже на реальной ошибке позы;
  \item совместная ветка держит покрытие (0{,}93--0{,}97), но отказывается в 60\,\% разбиений для велосипеда и в 86\,\% для растения. В остальных её радиус около 2~м;
  \item в миссиях (216 задач, у которых старт и цель остаются в пределах карты):
  \begin{itemize}[nosep]
    \item карта без калибровки находит путь лишь в 23\,\% задач, и $44\,\%$ этих путей нарушают безопасное расстояние, а $38\,\%$ проходят через истинные препятствия;
    \item ветка «только метки» находит путь в 8 задачах, 6 из них небезопасны;
    \item совместная ветка отказывается от миссии в 89\,\% задач.
  \end{itemize}
\end{itemize}
Иначе говоря, при одометрии без замыканий циклов полезного сертификата нет, и совместная калибровка честно об этом сообщает. Калибровка только по меткам выдаёт ложный сертификат. Чтобы гарантия стала полезной, нужна оценка позы с ограниченной ошибкой (SLAM с замыканиями, релокализация), а не другая калибровка.

\begin{figure}[H]
\centering
\includegraphics[width=\linewidth]{figures/variants.png}
\caption{Проверка упрощений: покрытие (сверху; линия --- номинал 0{,}9) и медианный радиус (снизу; штриховка --- ветка отказывается в половине разбиений или чаще) для масок MobileSAM, реальной визуальной одометрии и сцен HM3D}
\label{fig:variants}
\end{figure}

\textbf{Другие сцены (вне распределения).} Гарантия выведена для обменяемых сцен. Чтобы увидеть, что происходит без обменяемости, радиусы, откалиброванные на всех 21 сцене ReplicaCAD, проверены на сценах HM3D~\cite{hm3d} из OSMa-Bench. Это реальные сканы домов, опубликованные только в условии \texttt{no\_lights}.
\begin{itemize}[nosep]
  \item Категории HM3D сведены к классам ReplicaCAD (\texttt{plant}, \texttt{decorative plant} $\to$ растение, \texttt{bicycle} $\to$ велосипед), высота пола оценена по эталонным точкам пола.
  \item Из 8 сцен две многоэтажные исключены. Растение видно с траектории лишь в двух сценах из шести, велосипед и тумба --- ни в одной.
  \item В сцене \texttt{00813} растение найдено хорошо: промах 0{,}18~м, лучше медианы ReplicaCAD.
  \item В сцене \texttt{00815} декоративное растение детектор делит между \texttt{plant\_stand}, \texttt{sculpture} и \texttt{pot}, и промах 2{,}9~м.
\end{itemize}
Покрытие на двух сценах --- 0{,}5 при номинале 0{,}9 и у совместной, и у остальных веток. Две сцены --- это иллюстрация, а не оценка. Вывод совпадает с наблюдением~\cite{sundarsingh} о падении покрытия при сдвиге условий: гарантия не переносится между семействами сцен, калибровку нужно повторять в целевой среде.
